% s-048-p048-m002; original PDF 48, 49
\Aphorism{s-048-p048-m002}{39}
Menschen sind keine Principien und der verrechnet sich, der sie zu berechnen unternimmt. Für Einen, der mit der Natur und ihrem Gesetzmässigen zu thun hat, ist der Umgang mit Menschen eine Erholung, ein Sprung in’s Gegensätzliche, oft aber auch eine Geduldprobe: wir Verwöhnten, denen eine Musik in reinen Verhältnissen als geistiges Element genugthäte, finden, dass das wohltemperirte Klavier, genannt Mensch, ewig verstimmt und allzuleicht verstimmbar ist. Man gehe in die beste, vertrauteste Gesellschaft: Drei, Vier, die sich kennen, die Grund haben und ihre Humanität darin suchen, durch Schonung einander bei Laune zu erhalten — und doch, eines Tages will die Wärme, die sociale Temperatur sich nicht einstellen, ohne dass ein fremder Störenfried nachweisbar wäre. Man ist geneigt zu reden und zu hören wie sonst, man ist auf einander vorbereitet und hat keine Überraschung zu scheuen, jeder Einzelne ist auch im Reinen mit sich selbst (dies die erste Grundlage der guten Geselligkeit): alle Versuchsbedingungen sind erfüllt wie sonst und trotzdem entfernt sich der wirkliche Erfolg vom berechneten vielleicht so weit, dass die Freunde beim Auseinandergehen nichts weniger als ein Wiedersehen zu wünschen haben. Was geschehen ist, weiss Keiner von ihnen: vielleicht sassen Zwei nebeneinander, die sonst gegenüber zu sitzen pflegten, vielleicht wurde ein Wort mit einer zufälligen Betonung ausgesprochen, die eine lange verklungene Saite erzittern machte — wer will den feinen Schiebungen und Gleitungen nachspüren, durch die ein menschliches Verhältniss aus seiner Schwebe gerückt wird? Die zarten Dinge sind auch die flüchtigsten, und alles, was eine Relation zwischen zwei Seelen ausdrückt, ist auf eines Messers Schneide gestellt: es ist, in der Mechanik der Gefühle, der Fall des labilen Gleichgewichts, das durch den leisesten Anstoss dauernd gestört wird. Für den Verkehr mit Leuten genügt Vorsicht und eine gelenke Hand; der Umgang mit Menschen ist eine blosse Möglichkeit, ein Grenzfall, den keine Berechnung erzwingen, den nur ein gnädiger Augenblick verwirklichen kann.
\par

